% !PRINTED 352
\weberpage{352}
%<<CONT
is. We also assume that the curve of the fourth order is free of singular points\footnote{\gesperrt{Jacobi}, \glqq Ueber die Anzahl der Doppeltangenten ebener algebraischer Curven\grqq{}, \gesperrt{Crelle's} Journal, Bd.\,40 (1850). \gesperrt{Clebsch}, \glqq Bemerkung zu \gesperrt{Jacobi's} Beweis für die Anzahl der Doppeltangenten\grqq{}, ebend., Bd.\,63 (1864). From the rather extensive literature on the theory of the double tangents of a curve of the fourth order the following writings are to be singled out. First, several papers by \gesperrt{Hesse} in \gesperrt{Crelle's} Journal, Bd.\,41, 49, 52. Further \gesperrt{Steiner}, Eigenschaften der Curven vierter Ordnung rücksichtlich ihrer Doppeltangenten, \gesperrt{Crelle's} Journal, Bd.\,49. \gesperrt{Aronhold}, Monatsber. d. Berl. Akademie 1864. \gesperrt{Geiser}, \glqq Ueber die Doppeltangenten einer ebenen Curve 4\,ten Grades\grqq{}, Math. Annalen, Bd.\,1 (1868). \gesperrt{Cayley}, \gesperrt{Crelle's} Journal, Bd.\,68 (1868). The work of \gesperrt{Salmon} cited above is also to be singled out here. In more recent times the theory of the double tangents has been developed further in connection with the theory of the \gesperrt{Abelian} functions. \gesperrt{Riemann}, Gesammelte mathematische Werke (2.\,Aufl., 1892) XXXI, aus dem Nachlass. \gesperrt{Weber}, Theorie der \gesperrt{Abel'schen} Functionen vom Geschlecht 3, Berlin 1876. \gesperrt{Frobenius}, \gesperrt{Crelle's} Journal, Bd.\,99 u.\,103. The algebraic side of the question is treated in two memoirs: \gesperrt{Nöther}, Math. Annalen, Bd.\,15 and \gesperrt{Weber}, ebend., Bd.\,23.}. Now let $f(x_1, x_2, x_3)$, or more briefly $f(x)$, be a ternary form of degree 4 which, set equal to zero, represents a curve of the fourth order without singular point, which we call the curve $f$.

If in this equation
\weq{1}{f(x_1, x_2, x_3) = 0}
we put in place of the variables $x_1, x_2, x_3$ expressions of the form

\[
x_1+t y_1, \quad x_2+t y_2, \quad x_3+t y_3,
\]

then there results an equation of degree 4 with respect to $t$:
\weq{2}{f(x+ty) = 0,}
whose roots, if $(x)$ and $(y)$ are two fixed points, when substituted in $x+ty$, give the coordinates of the points of intersection of the connecting line of the points $(x), (y)$, which we shall call the line $(x,y)$, with the curve $f$.

Now let the function $f(x+ty)$ be ordered by powers of $t$. This gives (by Bd.\ I, \S.\,60):
\weq{3}{f(x+yt) = P_0(x,y)+4tP_1(x,y)+6t^2P_2(x,y)+4t^3P_3(x,y)+t^4P_4(x,y),}
in which the $P_\nu(x,y)$ are the polars of $f(x)$, and hence homogeneous
%CONT>>
